% Part: lambda-calculus
% Chapter: syntax
% Section: terms

\documentclass[../../../include/open-logic-section]{subfiles}

\begin{document}

\olfileid{lam}{syn}{trm}
\olsection{ٹرم}

لامبڈا کیلکولس دے ٹرم استقرا نال بنائے جاندے نیں۔ ایہناں لئی
لامتناہی متغیر $\Obj{v_0}$، $\Obj{v_1}$، \dots، نشان~«$\lambd$»،
تے قوساں ورتدے آں۔ اسیں $x$، $y$، $z$، \dots{} نوں متغیراں
لئی تے $M$، $N$، $P$، \dots{} نوں ٹرماں لئی ورتاں گے۔

\begin{defn}[ٹرم] \ollabel{defn:term}
لامبڈا کیلکولس دے \emph{ٹرماں} دا سیٹ استقرا نال ایویں متعین کردے آں:
\begin{enumerate}
  \item \ollabel{defn:term-var} جے $x$ متغیر اے، تاں $x$ اک ٹرم اے۔
  \item \ollabel{defn:term-abs} جے $x$ متغیر اے تے $M$ ٹرم اے،
    تاں $(\lambd[x][M])$ اک ٹرم اے۔
  \item \ollabel{defn:term-app} جے $M$ تے $N$ دوویں ٹرم نیں،
    تاں $(MN)$ اک ٹرم اے۔
\end{enumerate}
\end{defn}

جے ٹرم $(\lambd[x][M])$ قاعدے \olref{defn:term-abs} نال
بنیا اے، تاں اسیں آکھدے آں کہ ایہ \emph{تجرید} دا نتیجہ اے۔
$x$ جیہڑا $\lambd[x]$ وچ اے، اوہنوں \emph{!!{parameter}} آکھدے آں۔
ٹرم $(MN)$ جیہڑا قاعدے \olref{defn:term-app} نال بنیا اے،
\emph{اطلاق} دا نتیجہ اے۔

اُپر دِتّی تعریف وچ ٹرماں دیاں ساریاں قوساں لکھی جاندیاں نیں۔
ایہ لکھت بھاری ہو سکدی اے، جیویں ٹرم
$(\lambd[x][((\lambd[x][x])(\lambd[x][(xx)]))])$ توں دِسدا اے۔
اسیں قوساں چھڈن دے رواج اگے دساں گے۔ رسمی تعریف نال پر ایہ
لبھنا سوکھا اے کہ قاعدیاں \olref{defn:term} مطابق ٹرم کیویں بنیا۔
مثال لئی، ٹرم $(\lambd[x][((\lambd[x][x])(\lambd[x][(xx)]))])$
بناؤن دا آخری قدم لازماً تجرید اے، جیہدا !!{parameter}~$x$ اے۔
تجرید ٹرم $((\lambd[x][x])(\lambd[x][(xx)]))$ توں شروع ہوندی اے،
جیہڑا دو ٹرماں دا اطلاق اے۔ اوہناں وچوں ہر ٹرم آپ تجرید دا
نتیجہ اے۔ ایس طرح اسیں اگے ودھ سکدے آں۔

\begin{prob}
دسو کہ $(\lambd[g][(\lambd[x][(g (x x))])
  (\lambd[x][(g (x x))])])$ کیویں بنیا اے۔
\end{prob}

\end{document}
